\documentclass{article}

% Load the cheatsheet package and necessary math packages
\usepackage{cheatsheet}
\usepackage{amsmath,amsfonts,amssymb,bm}
\usepackage{graphicx}
\usepackage{verbatim}

% Custom macros for mathematical notation
\newcommand{\vb}[1]{\mathbf{#1}} % Vector (bold)
\newcommand{\mt}[1]{\mathit{#1}} % Matrix (italic)
\newcommand{\RR}{\mathbb{R}} % Real numbers
\newcommand{\vn}[1]{\left\|#1\right\|} % Magnitude/Norm
\newcommand{\vs}[1]{\mathcal{#1}} % Vector space
\newcommand{\proj}[2]{\text{Proj}_{#1}(#2)} % Projection of #2 onto #1
\DeclareMathOperator{\tr}{tr} % Trace operator

\newcommand*{\todo}{{\color{red}\textbf{TODO:}} }
\newcommand*{\bP}{\mathbf{P}}
\newcommand*{\bp}{\mathbf{p}}
\newcommand*{\bq}{\mathbf{q}}

% Customize cheatsheet parameters
\setcheatsheetcolumns{3}
\setcheatsheetfontsize{6.0}  % Tune font size for readability/compactness
\setcheatsheetmargin{0.1in}
\setcheatsheetdensity{0.05}
% \hidecheatsheetheader      % No header needed

\renewcommand{\studentname}{Nicholas Schlabach}
\renewcommand{\classname}{Math 468, Stochastic Processes}
\renewcommand{\sheetdate}{April 22nd, 2026}

\begin{document}

% Begin the cheat sheet
\begincheatsheet
\setcounter{seccounter}{-1}
\nsection{Facts}
\nsubsection{Algebra}{
\begin{itemize}
    \item Geometric Series: $\sum_{k=0}^n r^k=\frac{1-r^{n+1}}{1-r}$
\end{itemize}}



% Probability Review taken from the 466 cheatsheet
\nsection{Probability Review}
\nsubsection{Basics}{
\begin{itemize}
    \item If $g$ is increasing and $Y=g(X)$, then $F_Y(y)=F_X(g^{-1}(y))$ and $f_Y(y)=f_X(g^{-1}(y))\cdot (g^{-1})'(y)$.
    \item $V[X]=E[X^2]-E[X]^2=E\big[(X-E[X])^2\big]$
    \item \textit{Def}: $Cov(X,Y)=E[XY]-E[X]E[Y]$. It follows that $V[X\pm Y]=V[X]+V[Y]\pm 2\text{Cov}(X,Y)$.
    \item $V[aX-b]=a^2V[X]$
    \item LOTUS: $E[g(X)]=\int_{-\infty}^\infty g(x)f_X(x)dx$
    \item \textit{Def}: $F_X(x)=P(X \le x), \; F_{X,Y}(x,y)=P(X\le x, Y \le y) \implies \frac{\partial^2}{\partial x \partial y}F_{X,Y}(x,y)=f_{X,Y}(x,y)$
    \item Conditional density is $f_{Y|X}(y|x)=\frac{f_{X,Y}(x,y)}{f_X(x)}$.
\end{itemize}
}
\nsubsection{Moment Generating Functions}{
\begin{itemize}
    \item \textit{Def}: $M_X(t)=E[e^{tX}] \implies E[X^n]=M^{(n)}_X(0)$
\end{itemize}
}
\nsubsection{Conditional Probability \& Independence}{
\begin{itemize}
    \item \textit{Def}: Events are independent if $P(A\cap B)=P(A)P(B)$. RVs are independent if the PDF/PMF factors.
    \item $P(A|B)=P(A \cap B)/P(B)$
    \item Law of Total Probability: If $B_j$ is a partition of $\Omega$ then $P(A)=\sum P(A|B_j)P(B_j)$
    \item If $A,B$ are independent all of the following occur:
        \begin{itemize}
        \item $E[AB]=E[A]E[B], \; Cov(A,B)=0$,
        \item $g(A),g(B)$ are independent regardless of $g,f: \mathbb{R} \to \mathbb{R}$
        \item $M_{X+Y}=M_X(t)M_Y(t)$ (Ch 7, Thm. 6)
        \end{itemize}
\end{itemize}
}

\nsubsection{Facts about RVs}{
\begin{itemize}
    \item Both exponential and geometric RVs have the loss of memory property $P(X>n+m|X>m)=P(X>n)$
\end{itemize}}

\nsection{Chapter 2}
\nsubsection{Definitions}{
A random (or stochastic process) is a family of random variables on $X_t$ so that $X_t: \Omega\to S$ is a function on a probability space $\Omega$. For midterm 1, time and $S$ are both discrete. \\
We say that $X$ is a \textit{Markov chain} if $P(X_{n+1}=x_{n+1}|X_0=x_0,X_1=x_1,\dots,X_n=x_n)=P(X_{n+1}=x_{n+1}|X_n=x_n)$. We also assume that a Markov chain is homogenous, meaning that the transition probability is also independent of time $n$. \\
To be a probability matrix, \textit{rows} must sum to one and entries need to be nonnegative. \\
\begin{itemize}
    \item $T_A=\min\{n{>}0:X_n\in A\}$ is  a RV called "hitting time"
    \item $N(y)=\#\{X_i=y:i\in \mathbb{Z}^+\}$ is the RV number of hits
    \item $N_n(y)$ is the number of hits up to time $n$
    \item $\rho_{xy}:=P_x(T_y<\infty)$ is probability $x$ reaches $y$
    \item $x \in S$ is an absorbing state if $\bP(x,x)=1$
    \item $y \in S$ is recurrent if $P_y(T_y<\infty)=1$, is transient otherwise
    \item $x$ leads to $y$ ($x \to y$) means $\rho_{xy}>0$
    \item $x$ and $y$ communicate  $(x\leftrightarrow y$) means $x\to y$ and $y\to x$
    \item A set $C\subseteq S$ of the form $\{y:y\leftrightarrow x\}$ for some fixed $x$ is called a communication class
    \item A Markov chain is irreducible if $S$ is a single communication class
    \item A set $C \subseteq S$ is closed if $\forall x \in C,x\to y \implies y\in C$
    \item An equilibrium distribution $\pi$ satisfies $\pi\bP=\pi$
    \item This equilibrium is stable iff $\lim_{n \to \infty}\bP^n(x,y)=\pi(y)$ for every $x,y \in S$. Equivalently, for every initial distribution $\pi_0$, $\pi_0\bP^n\to \pi$.
    \item $m_y:=E_y(T_y)$ is the expected return time. A recurrent state is called null-recurrent if  $m_y=\infty$ and positive recurrent otherwise
\end{itemize}
}
\nsubsection{Basic Properties}{
\textbf{Thm 2.1}: If $y$ is a transient state $(\rho_yy<1)$, then no matter where the chain starts the chain visits $y$ a finite number of times and the expected number of visits is finite. For any state $x$, $P_x(N(y)<\infty)=1$ and $E_x(N(y))=\sum_{n=1}^\infty \bP^n(x,y)<\infty$. \\
On other hand, if $y$ is recurrent and starts at $y$ than the chain visits $y$ an  infinite number of times $P_y(N(y)=\infty)=1$ and $E_y(N(y))=\infty$. More generally $P_x(N(y)=\infty)=\rho_{xy}$, $E_x(N(y))=0$ if $\rho_{xy}=0$ and $E_x(N(y))=\infty$ if $\rho_{xy}>0$
\textbf{Corollary 1}: A state $x$ is transient iff $E_X(N(x))<\infty$ \\
\textbf{Thm 2.2}: (a) If $x \in S$ is recurrent and $x\to y$, then $Y$ is recurrent, $\rho_{xy}=\rho_{yx}=1$ and $y \to x$. Hence, recurrence is shared across communication classes. (b) The same is true for null recurrent and positive recurrent states. (c) If $S$ is finite, then there is at least one recurrent state and recurrent states are positive recurrent. (d) The set of recurrent states $S_R$ is the disjoint union of communication classes.
}
\nsubsection{Equilibrium}{
The equilibrium vanishes on transient states. \\
\textbf{Thm 2.3}: Let $N_n(y)$ be the number of visits to $y$ in time $1\le t \le n$, then $\lim_{n\to\infty} \frac{N_N(y)}{n}=\frac{I_{\{T_y<\infty\}}}{m_y}$ with probability one, where $I$ is the indicator function. \\
\textbf{Thm 2.4}: (a) An equilibrium $\pi$ vanishes on all transient and null recurrent states. (b) A positive recurrent communication class supports a unique equilibrium $\pi(y)=\frac{1}{m_y}$ for $m_y$ finite. (c) Finite closed communication class is automatically positive recurrent
}
\nsubsection{Stability of Equilibrum}{
Define the period as $d_x=gcd\{n:\bP^n(x,x)>0\}$. \\
Periods are the same for nodes that communicate. \\
\textbf{Thm 2.5:} If $C$ is a positive recurrent class and $\pi$ is the unique equilibrium supported on $C$.
(a) If $C$ is aperiodic, then $\pi$ is a stable equilbirum \\
The Peron-Frobenus thm says real square matrix with positive entries there is a unique eigenvalue of positive magnitude, that eigenvalue is real, has multiplity one, and the eigenvector can be chosen to have nonnegative entries. \\
\textbf{Them 2.6}: If a chain is finite, irreducible (hence positive recurrent) and for some power $m$, all entreis $P^m$ are nonzero, then the chain is aperiodic and the unique equilibrum ist stable.
}
\nsubsection{Misceleaneous}{Note that the formula $\pi(y)=\frac{1}{y}$ for finite states works both ways, if you know the equilbirum then you know the expected return times. \\
The following may be useful:
\[ E_X(N(y))=\sum_{n=1}^\infty \bP^n(x,y) \]
There is another thm stating that if $y$ is transient then the chain visits any other state $y$ a finite number of times and $E_x(N(y))=\frac{\rho_{xy}}{1-\rho_{yy}}$ \\
The following is useful:
\[\bP^n(x,y) =\sum_{m=1}^n P_x(T_y=m)\bP^{n-m}(y,y) \]
}
\nsection{Markov Jump Processes}
\nsubsection{Basics}{

Let $X(t)$ be a stochastic process with discrete space $S$. The RVVs $\tau_1,\tau_2\dots$ denote the jumps for the process. We say $X(t)$ is a Markov jump process if it satisfies the markov property $P(X(t)=y|X(s_1)=x_1\dots,X(s_n)=x_n,X(s)=x)=\bP(t-s,x,y)$
\begin{itemize}
    \item For a Markov jump process $P(X(0){=}x_0,\dots,X(t_n){=}x_n)=\pi(x_0)\bp(t_1;x_0,x_1)\dots \bp(t_n-t_n{-1},;x_{n-1},x_n)$
    \item The transition function satisfies the Chapman-Kolmogorov equation $\bP(s+t)=\bP(s)\bP(t)$ or $\bP(s+t);x,z)=\sum_{y \in S}\BP(s;x,y)\bP(t;y,z)$. Then $\bP(t)=e^{t\bq}$ where $\bq=\bp'(0)$. 
    \item For $x\ne y$, $b\q(x,y)$ is the rate at which the chain jumps from $x$ to $y$ and $q_x=-\bq(x,x)=\sum_{y \ne x}\bq(x,y)$ is the rate at which the chain jumps from $x$ to $y$.
    \item If $\tau_1,\tau_2,\dots$ denote the times of the jumps, then $\tau_1,\tau_2-\tau_1,\dots$ are independent random variables. If the chain starts in state $x$, then $\tau_1$ is an exponential RV with parameter $q_x$
\end{itemize}}
\nsubsection{Poisson Process}{
\begin{itemize}
    \item In the case of a Poisson process, the states are $\Z$ where $\bq$ has entries $-\lambda$ on the diagonal and $\lambda$ for each entry directly above the diagonal.
    \item The following characterizes the Poisson process:
        \begin{itemize}
        \item $X(0)=0$ and $X(t_1),X(t_2)-X(t_2),\dots$ are independent
        \item $X(t)=x(s)$ has Poisson distribution with parameter $\lambda(t-s)$ 
        \end{itemize}
    \item
\end{itemize}
}
\nsubsection{Important Properties}{
\begin{itemize}
    \item For a stochastic matrix $\bP$, $\lambda=1$ is an eigenvalue and the other eigenvalues lie in $\mathbb{C}$ and hae magnitude less than one. A Markov jump process is determined by an initial dist. and the infinitesimal generator matrix $\bq(x,y)$. Given $\bq$, the dynamics of the process are dtermiend by the family of stocahstic matrices $\bP(t)=e^{t\bq}$ satisyfing Chapman-Kolmogorov $\bP(t+s)=\bP(t)\bP(s)$.
    \item Given a Markiov jump process
        \begin{enumerate}
        \item The off-diagonal entries of $\bq$ are nonnegative, and the diagonal entries are negative. The sum of the entries on any row is zero.
        \item It follows zero is an eigenvalue with eigenvector $\mathbf{1}$ and the other eigenvalues have negative real part (this follows from the similar stmt. on $\bP$)
        \end{enumerate}
        Conversely, given a matrix $\bq$ satisfying (1) there exists a Markov jump process with finite state space having $\bq$ from a generator. Part (a) follows from expanding $P(t)=1+tq+o(t)$
    \item
        We define $T_y\min\{t \ge \tau_1:X(t)=y\}$ and $\rho_{x,y}=P_x(T_y<\infty$) similar to Markov chains. Say a process is recurrent/transient if all states have that property individually.
    \item If $x$ is nonabsorbing, then $Q(x,Y)=P_x(X(\tau_1)=y$. If $x$ is absorbing, then $\overset{\sim}{P}(x,x)=1$ and for $\overset{\sim}{P}(x,y)=0$ when $x\ne y$.
    \item An equilbrium is a dist. such that $\piP(t)=\pi$.
    \item Then $Q$ is a probability transition matrix for a Markov chain which is called the embedded chain (roughly this is $X(\tau_1)X(\tau_2),\dots$).
    \item The quantities $\rho_{x,y}$ for Makov jump processes are the same as the corresponding quanitties for the embedded chain. Then $\pi(y)=\frac{1}{q_ym_y}$ where $m_y=E_Y(T_y)$ is the expected return time.
    \item $q(x,y)=q_xQ(x,y)$. In words, the rate of change for the system to leave $x$ (which will occur at time $\tau_1)$ and go to $y$ is equal to the rate of change for the system to leave $x$ times the probability that the system will go from $x$ to $y$
\end{itemize}
}
\nsection{Brownian Motion \& Stochastic Integrals}
\nsubsection{Brownian Motion}{By def, Brownian motion (or the Wiener process) satisfies
\begin{enumerate}
    \item $W(0)=0$ and $W(t)-W(s)\sim N(0,t-s)$ is a normal RV.
    \item $W(t_1)-W(s_1),W(t_2)-W(s_2),\dots W(t_n)-W(s_n)$ are independent for $0\le s_1\le t_1\le s_2\le t_2 \dots \le s_n\le t_n$
\end{enumerate}
It follows from this definition that $E[(W(t)-W(s))^2/(t-s)=1$
The following properties are also true:
\begin{itemize}
\item $W(t)$ is continuous with probability 1
\item $W(t)$ is not differentiable with probability one
\item Given any interval $I$, $W(t)$ has infinitely many zeros in $I$, making it recurrent.    
\end{itemize}}
\nsubsection{Stochastic Integrals}{
We define the stochastic integral of $f$ relative to Brownian motion as
\[\int_a^b f(t)dW(t)=f(b)W(b)-f(a)W(a)-\int_a^bf'(t)W(t)dt \]
We call $dW(t)$ white noise. This integral has the following properties
\begin{itemize}
    \item The integral is linear in $f$
    \item $E[\int_a^bf(t)dW(t)]=0, V[\int_a^bf(t)dW(t)]=\int_a^b f(t)^2dt$. Intuitively this follows from $dW(t)$ being "independent"
    \item Product of expectations is expectation of product
        \[ E\Big[\Big(\int_a^bf(t)dW(t)\Big)\Big(\int_a^b g(t)dW(t)\Big)\Big]=\int_a^bf(t)g(t)dW(t)\ \]
\end{itemize}
}
\nsubsection{Ito Stochastic Integral}{
% NOTE: It matters how we do t_
We define the Ito stochastic integral using a Riemann sum $Z(t)=\int_0^tY(s)dW(s)=\lim \sum Y(t_k)(W(t_{k+1})-W(t_k))$  \\
This is linear in $Y$, has mean zero and is a cont. function of $t$. The variance is $V[Z(t)]=\int_0^t E[Y(s)^2)ds$.  \\
Now suppose that $F \in C^2$. Then $Y(t)=F(W(t))$ is a non-anticipating process compared to Brownian motion and 
\[ dF(W(t))=F'(W(t))dW(t)+\frac{1}{2}F''(W(t))dt \]
Integrating gives $\int_0^t dF(W(t))=F(W(t))-F(W(0))=\dots$ \\
Even more generally, with $Y(t)=F(t(W(t))$ we get  
$dF(t,W(t))=F'(t,W(t))dW(t)+\big(\dot{F}(t,W(t))+\frac{1}{2}F''(t,W(t))\big)dt $.
}

% End the cheat sheet
\endcheatsheet
\end{document}
